\section{Related Work}
\label{sec:related_work}

\subsection{Time Series Foundation Models}
Prior to foundation models, deep learning for time series relied heavily on dataset-specific architectures, evolving from RNNs \cite{deepar} and MLPs \cite{nbeats} to sophisticated Transformers \cite{informer,autoformer,patchtst,itransformer} and multi-scale mixers \cite{timesnet,timemixer}. Recently, the paradigm shifted toward zero-shot generalization. Large language models have been repurposed for forecasting \cite{llmtime,timellm}, while native time series models like Chronos \cite{chronos} tokenize time series into discrete bins and train a language model to predict subsequent tokens, while models like Moirai \cite{moirai} and TimesFM \cite{timesfm} employ patch-based attention mechanisms. Despite their zero-shot capabilities, these models are bounded by strict context window limits, restricting their ability to leverage long-term historical structural patterns.

\subsection{Retrieval-Augmented Time Series}
To ground foundation models in broader historical context, retrieval mechanisms have been proposed. TS-RAG \cite{tsrag} employs a neural encoder to fetch relevant historical sub-series, projecting them via a pre-trained Adaptive Retrieval Mixer (ARM). Similarly, TimeRAF \cite{timeraf} uses embedding similarity and neural fusion. While these methods achieve state-of-the-art accuracy on continuous benchmarks like ETT and Weather, they depend entirely on deep latent projections to circumvent numerical scale mismatch, necessitating computationally expensive retraining phases.

\subsection{Non-Stationarity and Statistical Retrieval}
The fragility of statistical similarity in non-stationary time series is widely recognized. Normalization techniques like Reversible Instance Normalization (RevIN) \cite{revin} attempt to mitigate distribution shifts internally within neural networks. On the retrieval front, SARAF \cite{saraf} addresses this by modulating retrieval diversity based on regime shifts, while TRACE \cite{trace} abandons raw numerical similarity in favor of text-based multimodal alignment. In contrast, our approach directly confronts the numerical brittleness of Euclidean k-NN. Rather than bypassing raw similarity via latent embeddings or semantic text, we introduce explicit mathematical scale restoration. This provides a deterministic, computationally lightweight solution to scale hallucination that remains entirely parameter-free.
